\documentclass[a4paper, 13pt, 3p, twoside]{report}
\usepackage{scrextend}
\changefontsizes{13pt}
\usepackage[utf8]{vietnam}
\usepackage[top=2cm, bottom=2cm, left=3.5cm, right=2.5cm]{geometry}
\usepackage{graphicx}
\usepackage{fancybox}
\usepackage{indentfirst}
\usepackage{amsthm}
\usepackage{latexsym}
\usepackage{amsmath}
\usepackage{amssymb}
\usepackage{times}
\usepackage{array}
\usepackage{enumitem}
\usepackage{subfiles}
\usepackage{titlesec}
\usepackage{titletoc}
\usepackage{chngcntr}
\usepackage{pdflscape}
\usepackage{afterpage}
\usepackage[ruled, vlined]{algorithm2e}
\usepackage{capt-of}
\usepackage{multirow}
\usepackage{fancyhdr}
\usepackage[natbib, backend=biber, style=ieee]{biblatex}
\usepackage{appendix}
\usepackage[font=small, labelfont=bf]{caption}
\usepackage{listings}
\usepackage{float}
\usepackage{subcaption}
\usepackage{xurl}
\usepackage[nonumberlist, nopostdot, nogroupskip, acronym]{glossaries}
\usepackage{glossary-superragged}
\setglossarystyle{superraggedheaderborder}
\usepackage{setspace}
\usepackage{parskip}
\usepackage{tocbasic}
\usepackage{blindtext}
\usepackage{indentfirst}

% ===================================================

\renewcommand{\bibname}{Danh_sach_tai_lieu_tham_khao} 
\usepackage[backend=biber,style=ieee]{biblatex}  %backend=biber is 'better'


\renewcommand\appendixname{PHỤ LỤC}
\renewcommand\appendixpagename{PHỤ LỤC}
\renewcommand\appendixtocname{PHỤ LỤC}


\addbibresource{Danh_sach_tai_lieu_tham_khao.bib} % chèn file chứa danh mục tài liệu tham khảo vào 

\include{lstlisting} % Phần này cho phép chèn code và formatting code như C, C++, Python

\input{Tu_viet_tat}
% ===================================================

\fancypagestyle{plain}{
    \fancyhf{}
    \fancyfoot[RO, RE]{\thepage}
    \renewcommand{\headrulewidth}{0pt}
    \renewcommand{\footnoterule}{0pt}
}

\setlength{\headheight}{10pt}

\def \TITLE{ĐỒ ÁN TỐT NGHIỆP}
\def \AUTHOR{LÊ ĐỨC LƯƠNG}

% ===================================================
\titleformat 
    {\chapter} % command
    [hang] % shape
    {\centering\bfseries} % format
    {CHƯƠNG \thechapter.\ } % label
    {0pt} %sep
    {} % before
    [] % after

\titlespacing{\section}{0pt}{\parskip}{0.5\parskip}

\titleformat
    {\subsection} % command
    [hang] % shape
    {\bfseries} % format
    {\thechapter.\arabic{section}.\arabic{subsection}\ \ \ \ } % label
    {0pt} %sep
    {} % before
    [] % after
\titlespacing{\subsection}{30pt}{\parskip}{0.5\parskip}

\renewcommand\thesubsubsection{\alph{subsubsection}}
\titleformat
    {\subsubsection} % command
    [hang] % shape
    {\bfseries} % format
    {\alph{subsubsection}, \ } % label
    {0pt} %sep
    {} % before
    [] % after
\titlespacing{\subsubsection}{50pt}{\parskip}{0.5\parskip}

% ===================================================

\usepackage{hyperref}
\hypersetup{pdfborder={0 0 0}}
\hypersetup{pdftitle={\TITLE},
    pdfauthor={\AUTHOR}}
\usepackage[all]{hypcap}
\graphicspath{{figures/}{../figures/}} % Thư mục chứa các hình ảnh
\counterwithin{figure}{chapter} % Đánh số hình ảnh kèm theo chapter. Ví dụ: Hình 1.1, 1.2,..
\title{\bf \TITLE}
\author{\AUTHOR}
\setcounter{secnumdepth}{3}
% \setcounter{tocdepth}{3}
\theoremstyle{definition}
\newtheorem{example}{Ví dụ}[chapter]
\onehalfspacing
\setlength{\parskip}{6pt}
\setlength{\parindent}{15pt}



% =========================== BODY ===============
\begin{document}
% \newgeometry{top=2cm, bottom=2cm, left=2cm, right=2cm}
\subfile{Bia} % Phần bìa
% \restoregeometry

% ===================================================
\pagenumbering{roman}
% \pagestyle{empty} % Header và footer rỗng
%\newpage
%\subfile{chapters/0_1_subject.tex}

\newpage
\pagenumbering{gobble}
\subfile{Chuong/0_2_Loi_cam_on.tex}

\newpage
\pagenumbering{gobble}
\subfile{Chuong/0_3_Tom_tat_noi_dung.tex}

% \newpage
% \pagenumbering{gobble}
% \subfile{Chuong/0_4_Tom_tat_noi_dung_English.tex}

% ===================================================
% \pagestyle{empty} % Header và footer rỗng
\newpage
\pagenumbering{roman} % Xóa page numbering ở cuối trang
\renewcommand*\contentsname{MỤC LỤC}

\titlecontents{chapter}
    [0.0cm]             % left margin
    {\bfseries\vspace{0.3cm}}                  % above code
    {{\bfseries{\scshape}
    CHƯƠNG \thecontentslabel.\ }}
    % numbered format
    {}         % unnumbered format
    {\titlerule*[0.3pc]{.}\contentspage}         % filler-page-format, e.g dots

    
\titlecontents{section}
    [0.0cm]             % left margin
    {\vspace{0.3cm}}                  % above code
    {\thecontentslabel \ } % numbered format
    {}         % unnumbered format
    {\titlerule*[0.3pc]{.}\contentspage}         % filler-page-format, e.g dots
    
\titlecontents{subsection}
    [1.0cm]             % left margin
    {\vspace{0.3cm}}                  % above code
    {\thecontentslabel \ } % numbered format
    {}         % unnumbered format
    {\titlerule*[0.3pc]{.}\contentspage}         % filler-page-format, e.g dots

 % Tạo mục lục tự động
\addtocontents{toc}{\protect\thispagestyle{empty}}
\tableofcontents 
\thispagestyle{empty}
\cleardoublepage

% \pagenumbering{roman}
%Tạo danh mục hình vẽ.
\renewcommand{\listfigurename}{DANH MỤC HÌNH VẼ}
{\let\oldnumberline\numberline
\renewcommand{\numberline}{Hình~\oldnumberline}
\listoffigures} 
% \phantomsection\addcontentsline{toc}{section}{\numberline {} DANH MỤC HÌNH VẼ}
\newpage


 %Tạo danh mục bảng biểu.
\renewcommand{\listtablename}{DANH MỤC BẢNG BIỂU}
{\let\oldnumberline\numberline
\renewcommand{\numberline}{Bảng~\oldnumberline}
\listoftables}
% \phantomsection\addcontentsline{toc}{section}{\numberline {} DANH MỤC BẢNG BIỂU}

\glsaddall 
 \renewcommand*{\glossaryname}{Danh sách thuật ngữ}
\renewcommand*{\acronymname}{DANH MỤC THUẬT NGỮ VÀ TỪ VIẾT TẮT}
\renewcommand*{\entryname}{Thuật ngữ}
\renewcommand*{\descriptionname}{Ý nghĩa}
\printnoidxglossaries
 %\phantomsection\addcontentsline{toc}{section}{\numberline {} DANH MỤC THUẬT NGỮ VÀ TỪ VIẾT TẮT}

% \newpage
 %\subfile{chapters/0_5_Danh_muc_viet_tat.tex}

% \newpage
% \subfile{chapters/0_6_Thuat_ngu.tex}



\newpage
\pagenumbering{arabic}

\pagestyle{fancy}
\fancyhf{}
\fancyhead[RE, LO]{\leftmark}
%\fancyhead[LE]{\rightmark}
\fancyfoot[RE, LO]{\thepage}

\chapter{GIỚI THIỆU ĐỀ TÀI}
\label{chapter:Introduction}
\subfile{Chuong/1_Gioi_thieu} % Phần mở đầu

\newpage
%\pagestyle{fancy} % Áp dụng header và footer
\chapter{KHẢO SÁT VÀ PHÂN TÍCH YÊU CẦU}
\label{chapter:Related_works}
\subfile{Chuong/2_Khao_sat}


\newpage
%\pagestyle{fancy} % Áp dụng header và footer
\chapter{CÔNG NGHỆ SỬ DỤNG}
\label{chapter:Methodology}
\subfile{Chuong/3_Cong_nghe}

\newpage
%\pagestyle{fancy} % Áp dụng header và footer
\chapter{THIẾT KẾ, TRIỂN KHAI VÀ ĐÁNH GIÁ HỆ THỐNG}
\label{chapter:Experiment}
\subfile{Chuong/4_Ket_qua_thuc_nghiem}

\newpage
%\pagestyle{fancy} % Áp dụng header và footer
\chapter{CÁC GIẢI PHÁP VÀ ĐÓNG GÓP NỔI BẬT}
\label{chapter:SolutionAndContribution}
\subfile{Chuong/5_Giai_phap_dong_gop}
\newpage
%\pagestyle{fancy} % Áp dụng header và footer
\chapter{KẾT LUẬN VÀ HƯỚNG PHÁT TRIỂN} %Kết luận và hướng phát triển}
\label{chapter:conclusion}
\subfile{Chuong/6_Ket_luan}

\newpage
\renewcommand\bibname{TÀI LIỆU THAM KHẢO}
\printbibliography
\phantomsection\addcontentsline{toc}{chapter}{TÀI LIỆU THAM KHẢO}

\appendixpage
\appendices
\addappheadtotoc
\end{document}